% 07-parser-expressions.tex — Chapter 7: Parser: Expressions
\chapter{Parser: Expressions}
\label{chap:07-parser-expressions}
\index{parser!expressions}

\pnum
Expression parsing is implemented in \srcref{src/parser/expr.h}.  The parser
uses a standard Pratt precedence-climbing algorithm for binary expressions,
with a separate recursive-descent layer for unary operators and primary
expressions.  The entry point is \cname{parse\_expr}.

% -------------------------------------------------------------------------
\section{Entry point and structure}
\label{sec:par-expr-entry}
% -------------------------------------------------------------------------

\pnum
\cname{parse\_expr(arena, parser)} captures the current source location, calls
\cname{parse\_binary\_expr(arena, parser, 0)}, and stamps the resulting node's
\cname{line}/\cname{col} fields.  The source location is captured before any
token is consumed, so it refers to the start of the expression.

\pnum
The call hierarchy is:

\begin{center}
\begin{tikzpicture}[node distance=0.7cm]
  \node[passbox] (e)   {\cname{parse\_expr}};
  \node[passbox, below=of e] (be)  {\cname{parse\_binary\_expr}};
  \node[passbox, below=of be] (ue) {\cname{parse\_unary\_expr}};
  \node[passbox, below=of ue] (pe) {\cname{parse\_primary\_expr}};
  \draw[arrow] (e) -- (be);
  \draw[arrow] (be) -- (ue);
  \draw[arrow] (ue) -- (pe);
\end{tikzpicture}
\end{center}

% -------------------------------------------------------------------------
\section{Pratt precedence climbing}
\label{sec:par-expr-pratt}
\index{Pratt parser}
\index{precedence climbing}
% -------------------------------------------------------------------------

\pnum
\cname{parse\_binary\_expr(arena, parser, min\_prec)}
(\srcref{src/parser/expr.h:27}) implements left-associative precedence climbing:

\begin{algorithm}
\textbf{parse\_binary\_expr}(\textit{arena, parser, min\_prec}):
\begin{enumerate}
\item Parse a unary expression as \textit{left}.
\item If the current token is \lain{as}, parse a cast type and wrap \textit{left}.
\item Loop:
  \begin{enumerate}[(a)]
  \item Let \textit{prec} = \cname{get\_precedence(token.kind)}.
  \item If \textit{prec} $<$ \textit{min\_prec}, break.
  \item Consume the operator token.
  \item Recurse: \textit{right} = \cname{parse\_binary\_expr(arena, parser, prec+1)}.
  \item Build \cname{EXPR\_BINARY(op, left, right)} and assign to \textit{left}.
  \item If the next token is \lain{as}, parse a cast and wrap \textit{left}.
  \end{enumerate}
\item Return \textit{left}.
\end{enumerate}
\end{algorithm}

\pnum
The recursive call with \textit{prec+1} produces left-associativity: operators
at the same precedence level associate to the left because the right operand
only claims operators of strictly higher precedence.

\pnum
The \lain{as} cast check after each sub-expression is required because \lain{as}
is a postfix operator with higher binding than binary operators.  The check after
the initial unary parse handles \lain{a as T op b}; the check inside the loop
handles \lain{a op b as T op c}.

% -------------------------------------------------------------------------
\section{Unary expressions}
\label{sec:par-expr-unary}
\index{unary expressions}
% -------------------------------------------------------------------------

\pnum
\cname{parse\_unary\_expr} (\srcref{src/parser/expr.h:60}) handles all prefix
unary operators:

\begin{center}
\begin{tabular}{lll}
\toprule
\textbf{Token} & \textbf{AST node} & \textbf{Meaning} \\
\midrule
\texttt{-}         & \cname{EXPR\_UNARY(TOKEN\_MINUS, right)}    & Arithmetic negation \\
\texttt{!}         & \cname{EXPR\_UNARY(TOKEN\_BANG, right)}     & Logical not \\
\texttt{\&}        & \cname{EXPR\_UNARY(TOKEN\_AMPERSAND, right)}& Address-of / borrow \\
\texttt{\~{}}      & \cname{EXPR\_UNARY(TOKEN\_TILDE, right)}    & Bitwise complement \\
\texttt{*}         & \cname{EXPR\_UNARY(TOKEN\_ASTERISK, right)} & Pointer dereference \\
\lain{mov}         & \cname{EXPR\_MOVE(right)}                   & Explicit ownership transfer \\
\lain{var}         & \cname{EXPR\_MUT(right)}                    & Mutable borrow \\
\bottomrule
\end{tabular}
\end{center}

\pnum
Unary operators are right-recursive: \cname{parse\_unary\_expr} calls itself
for the operand.  This allows chaining (\lain{--x} if supported, \lain{!not y},
etc.) and correctly handles \lain{mov *ptr}.

% -------------------------------------------------------------------------
\section{Primary expressions}
\label{sec:par-expr-primary}
\index{primary expressions}
% -------------------------------------------------------------------------

\pnum
\cname{parse\_primary\_expr} (\srcref{src/parser/expr.h:99}) handles all
atom-level and postfix constructs.

\subsection{Literals}

\pnum
Integer literals (\cname{TOKEN\_NUMBER}) are parsed with
\cname{parse\_numeric\_literal} and stored as \cname{EXPR\_LITERAL} with an
\texttt{int} value.  Float literals use \cname{strtod} and produce
\cname{EXPR\_FLOAT\_LITERAL}.  String literals produce \cname{EXPR\_STRING} with
the content pointer already stripped of delimiters (see \cref{sec:lex-strings}).
Character literals produce \cname{EXPR\_CHAR} with a decoded \texttt{char} value;
escape sequences (\texttt{\\n}, \texttt{\\t}, \texttt{\\\\}, \texttt{\\'}) are
decoded at parse time.

\subsection{Identifiers and member access}

\pnum
An \cname{TOKEN\_IDENTIFIER} produces an \cname{EXPR\_IDENTIFIER} node.  The
parser then checks for postfix operators in a loop:
\begin{itemize}
\item \texttt{.} followed by an identifier: wraps in \cname{EXPR\_MEMBER}.
\item \texttt{[} followed by an expression: produces \cname{EXPR\_INDEX}.
  If the subscript contains \texttt{..} or \texttt{..=}, the index expression
  is an \cname{EXPR\_RANGE}.
\item \texttt{(} followed by arguments: produces \cname{EXPR\_CALL}.
\end{itemize}

\pnum
Postfix operators are left-associative and bind more tightly than any binary
operator, so \lain{a.b[i](x)} parses as \lain{((a.b)[i])(x)}.

\subsection{Case expressions}

\pnum
A \lain{case} keyword inside an expression context initiates a match expression.
The optional \texttt{\&} prefix marks the scrutinee as borrowed
(\cname{is\_borrowed = true}).  Arms are parsed as \texttt{pattern => body}
pairs inside a brace-delimited block.  The \texttt{else} keyword introduces the
wildcard arm.

\subsection{Builtin values}

\pnum
\texttt{@os} and \texttt{@arch} produce \cname{EXPR\_BUILTIN} nodes carrying
the appropriate \cname{BuiltinKind} discriminant.  These are evaluated at
comptime by \cname{sema\_eval\_comptime\_expr()} in \srcref{src/sema/comptime.h}.

\subsection{Parenthesised expressions}

\pnum
A parenthesised expression \lain{(expr)} consumes \texttt{(}, calls
\cname{parse\_expr} recursively, then expects \texttt{)}.  The result is the
inner expression with no wrapper node.

\subsection{Array literals}

\pnum
\texttt{[} initiates an array literal.  Elements are comma-separated and
delimited by \texttt{]}.  The result is \cname{EXPR\_ARRAY\_LITERAL} with an
\cname{ExprList} of elements.

% -------------------------------------------------------------------------
\section{Cast expressions}
\label{sec:par-expr-cast}
\index{cast expressions}
% -------------------------------------------------------------------------

\pnum
The \lain{as} keyword desugars to \cname{EXPR\_CAST(expr, target\_type)}.  It is
treated as a postfix operator rather than a binary operator; this means it does
not participate in the standard precedence table and is checked explicitly after
each sub-expression during Pratt climbing.  The effect is that \lain{as} binds
tighter than all binary operators: \lain{a + b as T} parses as
\lain{a + (b as T)}.

\begin{rationale}
Treating \lain{as} as postfix prevents the ambiguity that arises if it were a
binary operator with a precedence level: \lain{a + b as T + c} would be
ambiguous as to whether \texttt{T + c} is a binary expression or part of the
type.
\end{rationale}
